% 本文件由 LaTeX 排版 Agent 根据有效 translation.json 自动生成。
% author=Hippolytus; work_id=0504; title=论世界的终结

\chapter{论世界的终结}

% structure: p0001 source=h1 target=omitted reason=page-title-already-rendered-by-parent

\Emph{真福殉道者、主教希玻律托的讲道，论世界的终结、假基督，以及吾主耶稣基督的第二次来临。}\par

% structure: p0003 source=h2 target=section reason=relative-heading-rank; base=h2

\section{1}

既然有福的先知们做了我们的眼睛，为我们的益处陈明了隐秘之事的清晰宣告，既通过生活，又通过宣讲，并通过圣神的默感，也论及尚未发生之事，这样他们也向世世代代描绘了最伟大的沉思与行动的主题。同样，他们也宣讲了天主在肉躯内降临世界，由无玷而承载天主的玛利亚以诞生和成长的方式降临，以及祂与众人相处的生命与言谈方式，祂藉圣洗的显现，那将要赐予众人的新生，藉圣洗的更新；还有祂的无数奇迹，祂在十字架上的蒙福苦难，祂从犹太人手中所受的凌辱，祂的埋葬，祂下到阴府，祂的再次上升，对古时灵魂的救赎，死亡的毁灭，祂从死者中赐予生命的复活，祂对全世界的再造，祂的升天与重返天堂，祂对圣神的接纳（宗徒们堪当领受），以及那将要宣告万有的第二次来临。因他们被称作\Emph{先知}，他们便必然预先指明并谈论了这些事。\par

% structure: p0005 source=h2 target=section reason=relative-heading-rank; base=h2

\section{2}

因此，他们也向我们指出了终结的日子，并预先表明了那将在末时出现、欺骗众人的背道者的日子，以及他王国的起始与终结，审判者的来临，义人的生命，罪人的惩罚，好叫我们众人，作为教会的子女，日日夜夜将这些事铭记于心，知道\Quote{这些事的一撇一捺也不会过去}\BibleRef{玛 5:18}，正如救主亲口所言。你们众人因此必须睁开你们心灵的眼睛和灵魂的耳朵，领受我们将要讲的话。因为今天我要向你们展开一个充满恐怖和恐惧的叙述，即关于终结的叙述，尤其是关于仇敌魔鬼对全世界的诱惑；此后，便是吾主耶稣基督的第二次来临。\par

% structure: p0007 source=h2 target=section reason=relative-heading-rank; base=h2

\section{3}

那么，基督的朋友们，我从何开始？以何开端？又当解说何事？对于所讲的事，我要提出什么证言？但让我们取那些（即先知们）作为我们开始这篇讲论的对象，并把他们作为可信的证人，以确认我们对所讨论之事所作的阐释；此后，还有宗徒们的教导，或者更恰当地说，预言，（为要明白）他们如何向全世界宣告终结之日的到来。既然这些先知们也已预先显明了尚未发生的事，并宣告了将要显现的恶人的诡计和欺骗，那么，来吧，让我们首先提出\Person{依撒意亚}作为我们的证人，因为他教导我们关于终结之时的事。那么，他说什么？\BibleQuote{你们的土地荒凉，你们的城市被火焚毁；你们的田地在你们眼前被外邦人侵吞；熙雍女子要像葡萄园里的茅舍，像黄瓜田中的草棚，像被围困的城邑。}\BibleRef{依 1:7}亲爱的，你们看见了先知的启示，他在如此多世代之前就宣告了那个时期。因为这话古时并非指犹太人，也不是指熙雍城，而是指教会。因为所有先知都宣告熙雍是从外邦人中迎娶的新娘。\par

% structure: p0009 source=h2 target=section reason=relative-heading-rank; base=h2

\section{4}

为此，让我们将话题转向第二个证人。这位证人如何？请听\Person{欧瑟亚}如此庄严地说话：\BibleQuote{在那些日子，上主要从旷野给他们带来一阵热风，使他们的血管干涸，使他们的泉源荒废；他们一切宝贵的器皿都要被掠夺。因为他们起来反对天主，必倒在刀下，他们的孕妇必被剖开。}\BibleRef{欧 13:15}这从东方而来的热风，除了那将要毁灭、使水和树木的果实干涸的假基督，还能是什么？因为那时人们全心投入他的工作。为此，他确实要毁灭他们，他们要在他的污秽中服事他。\par

% structure: p0011 source=h2 target=section reason=relative-heading-rank; base=h2

\section{5}

请注意先知与先知的一致。你也要认识另一位表达类似的先知。因为\Person{亚毛斯}以十分一致的方式预言了同样的事：\BibleQuote{上主这样说：因为你们用拳头殴打穷人，又夺取他们精选的礼物，你们建造房屋，却不得住在其中；栽种美好的葡萄园，却不得喝其中的酒。因为我深知你们的罪过众多，你们践踏正义，收受贿赂，在城门口屈枉穷人。因此，在那时候，明智人必缄默不语，因为那是邪恶的时代。}亲爱的，要知道那时人的邪恶，他们抢夺房屋田地，甚至剥夺义人的正义；因为这些事发生时，你们就知道这是终结。为此，你们藉着先知的智慧，以及那些日子将要显示的启示而受教。正如我们已说过的，所有先知都清楚指明了末时将要发生的事，正如他们也宣告了古时的事。\par

% structure: p0013 source=h2 target=section reason=relative-heading-rank; base=h2

\section{6}

\Quote{主论及那些使我子民迷途的先知这样说：他们用牙齿嚼食，却喊说\InnerQuote{和平}；若没有把食物放入他们口中，他们就预备与他争战。因此，为你们只有黑夜，没有神视；只有黑暗，不能占卜；太阳必不没落于先知之上，白昼在他们之上变为黑暗。先见者必蒙羞，占卜者必受辱。}\BibleRef{米 3:5}我们预先列举了这些事，为使你们知道末世要有的痛苦和混乱，以及众人彼此之间的生活方式、他们的嫉妒、仇恨、纷争，牧人对羊群的疏忽，百姓对司铎的悖逆态度。\par

% structure: p0015 source=h2 target=section reason=relative-heading-rank; base=h2

\section{7}

因此，众人都将随从自己的私意而行。儿女将向父母下手。妻子将把丈夫交于死地，丈夫将带自己的妻子如罪犯受审。主人将凶暴地压制仆人，仆人将向主人摆出悖逆的态度。无人会尊敬老人的白发，无人会怜悯少年的俊美。天主的圣殿将如同房屋，教堂处处被倾覆。圣经将被藐视，到处传唱仇敌的歌曲。奸淫、通奸、伪誓将充满大地；邪术、咒语、占卜将随着这一切，全力而热心地推行。总之，在那些自称为基督徒的人中，将兴起假先知、假宗徒、骗子、惹事生非者、作恶者、彼此撒谎者、奸淫者、通奸者、强盗、贪婪者、伪誓者、说谎者、互相仇恨者。牧人将像豺狼；司铎将拥抱谎言；隐修士将贪恋世俗之物；富人将心硬；统治者将不帮助穷人；有权势者将丢弃一切怜悯；审判官将夺去义人应得的正义，因受贿而眼瞎，他们将召来不义。\par

% structure: p0017 source=h2 target=section reason=relative-heading-rank; base=h2

\section{8}

关于人，我还说什么呢？既然连元素本身都将违背自己的秩序。各城将有地震，各地将有瘟疫；可怕的雷声和骇人的闪电将烧毁房屋和田野。暴风将过度地搅扰海陆；地上将有不结果实，海中有咆哮，因灵魂和人的毁灭而有难以忍受的骚动。太阳将有异兆，月亮将有异兆，星辰将有偏移，万民将有困苦，大气将有失常，冰雹将降于地面，冬季极其严寒，各种霜冻，无情的热风，意外的雷声，料想不到的火灾；总之，全地都有哀哭和悲伤，毫无安慰。因为，\Quote{由于罪恶的增加，许多人的爱情必要冷淡。}\BibleRef{玛 24:12}由于这一切的骚动和混乱，宇宙的主在福音中呼喊说：\Quote{你们要谨慎，不要受迷惑；因为有许多人要冒我的名来，说：我是基督，时候近了；所以不要跟从他们。但当你们听见战争和骚动时，不要惊慌；因为这些事必须先发生，但结局不会立刻来到。}\BibleRef{路 21:8}让我们留意救主的话，祂总是为我们的安全而告诫我们：\Quote{你们要谨慎，不要受迷惑；因为有许多人要冒我的名来，说：我是基督。}\par

% structure: p0019 source=h2 target=section reason=relative-heading-rank; base=h2

\section{9}

当祂被接升天回到圣父那里后，就有人起来，说：\Quote{我是基督，}如\Person{西满·术士}以及其余的人，他们的名字我们目前无暇提及。因此，在末后完成之日，必定会再有假基督起来，说：\Quote{我是基督，}他们要迷惑许多人。成群的人将从东奔到西，从北奔到海，说：基督在哪里？基督在那里？但他们怀着虚妄的念头，没有仔细阅读圣经，心术不正，他们将寻找一个找不到的名字。因为这些事必须先发生，然后丧亡之子——也就是说，魔鬼——必须出现。\par

% structure: p0021 source=h2 target=section reason=relative-heading-rank; base=h2

\section{10}

首先，\Person{伯多禄}，信德的磐石，基督，我们的天主，称为有福者，教会的导师，第一位门徒，持天国钥匙者，教导我们说：\Quote{第一要知道，孩子们，末世将有讥诮者，随从自己的私欲而行。在你们中间也必有假教师，私自引进可咒骂的异端。}随后，\Person{若望}神学家，基督的爱徒，与他一致地呼喊：\Quote{魔鬼的儿女是显而易见的；\BibleRef{若一 3:10} 现在已经有许多假基督；\BibleRef{若一 2:18} 但不要跟从他们。\BibleRef{路 21:8} 不可信每一个神，因为有许多假先知已经进入世界。\BibleRef{若一 4:1}}接着，\Person{犹达}，\Person{雅各伯}的兄弟，也这样说：\Quote{在末世将有讥诮者，随从自己不虔敬的私欲而行。这些人无所畏惧地喂养自己。}你们已看到了神学家和宗徒们的共识，以及他们教义的和谐。\par

% structure: p0023 source=h2 target=section reason=relative-heading-rank; base=h2

\section{11}

最后，请听\Person{保禄}如何大胆发言，并留意他多么清楚地揭示了这些事：\Quote{你们要提防那些作恶的工人，提防那些\OriginalTerm{割损}{\GreekText{κατατομή}}。\BibleRef{斐 3:2} 你们要小心，免得有人用哲学和虚妄的欺骗俘掠你们。\BibleRef{哥 2:8} 所以你们要谨慎行事，因为时日邪恶。}\BibleRef{弗 5:15} 总之，一个人若在教会中从先知、宗徒并从主自己那里听到了这些事，却仍不留心照顾自己的灵魂、不顾及终结的时候和那即将来到的、我们必须站在基督审判台前的时刻，他还有什么借口呢？\par

% structure: p0025 source=h2 target=section reason=relative-heading-rank; base=h2

\section{12}

但既然我们已完成了对终结的论述，我们将把讲解转向接下来要陈述的事。为此，我引证一位完全值得信赖的见证人——即先知\Person{达尼尔}，他解释了\Person{拿步高}的神视，并从诸王的开端直到其结局，为那些寻求行走在其中的人指明了正道——即真理的显示。因为先知说了什么？他以如下方式向\Person{拿步高}预先指明了这事：\Quote{王啊，你观看，见有一个大像立在你面前，这像的头是精金的，胸膛和手臂是银的，腹部和大腿是铜的，腿是铁的，脚是半铁半泥的。你观看，见有一块非人手凿出来的石头，打在这像半铁半泥的脚上，把脚砸碎；于是铁、泥、铜、银、金都一同砸得粉碎，如同夏天禾场上的糠秕；打碎这像的石头变成一座大山，充满全地。}\BibleRef{达 2:31}\par

% structure: p0027 source=h2 target=section reason=relative-heading-rank; base=h2

\section{13}

因此，我们将\Person{达尼尔}的神视与这些结合起来，将两者合并为一个叙述，并表明先知在神视中所见的是何等真实并与\Person{拿步高}预先所见的一致。因为先知这样说：\Quote{我达尼尔观看，见天上有四股大风，刮在大海之上。有四只大兽从海中上来，各不相同。头一只像狮子，有鹰的翅膀；我观看，直到它的翅膀被拔去，它从地上被举起来，用两脚站立像人一样，又给了它人心。又有一兽如熊，就是第二兽，它直立一边，口齿间衔着三根肋骨；有人吩咐它说：起来，吞吃多肉。此后我观看，又有一兽如豹，背上有四个鸟翅；这兽有四个头。此后我观看，见第四兽甚是可怕，极其强大，有大铁牙和铜爪，吞吃嚼碎，所余剩的用脚践踏；这兽与前三兽大不相同，它有十角。我观看这些角，见其中又长起一个小角，在这小角前面，有三角被连根拔起；看哪，这角中有眼像人的眼，有口说夸大的话。}\BibleRef{达 7:2}\par

% structure: p0029 source=h2 target=section reason=relative-heading-rank; base=h2

\section{14}

既然先知以奥秘方式所说的这些事似乎对所有人都难以理解，我们将不向心智健全的人隐瞒任何一点。先知提到第一只兽，即从海中上来的\Emph{狮子}，\Person{达尼尔}是指\Place{巴比伦}王国在世上建立；这也就是这像的\Quote{金头}。而说到它的\Quote{翅膀如鹰}，他表明\Person{拿步高}王被高举，并自高反对天主。然后他说它的\Quote{翅膀被拔去}，意思是他的荣耀被推翻：因为他被赶出王国。而说到\Quote{给了它人心，并让它像人一样用脚站立}，意思是它悔改了，并承认自己不过是人，而归光荣于天主。看，我已经如此解释了第一只兽的相似。\par

% structure: p0031 source=h2 target=section reason=relative-heading-rank; base=h2

\section{15}

然后，在狮子之后，先知看见第二只兽像熊，这代表波斯人；因为巴比伦之后，波斯人获得了主权。而说到\Quote{我看见它口中有三根肋骨}，他指的是三个民族：波斯人、玛待人和\Place{巴比伦}人，这些也由像中金之后的银所表达。看，我们也已经解释了第二只兽。然后第三只是豹，这代表\Place{希腊}人。因为在波斯之后，\Place{马其顿}王\Person{亚历山大}在消灭了\Person{达理阿}之后获得了主权；这由像中的铜所表达。而说到\Quote{鸟的四个翅膀，兽有四个头}，他非常清楚地表明\Person{亚历山大}的王国如何被分为四部分。因为它有四个头——即由它兴起的四个王。因为\Person{亚历山大}在临终时将自己的王国分为四部分。看，我们也已经讨论了第三只兽。\par

% structure: p0033 source=h2 target=section reason=relative-heading-rank; base=h2

\section{16}

接下来，他告诉我们关于\Quote{第四只兽，非常可怕而恐怖；它的牙齿是铁的，它的爪是铜的。}这些所指的不正是罗马王国吗？这王国也由铁所象征，它要粉碎它之前的所有帝国宝座，并统治整个大地。此后，先知所见的一切中，还有什么留给我们解释呢？无非是\Quote{像的脚趾，其中一部分是铁，一部分是泥，混杂在一起。}因为像的十个脚趾，他形象地指代从中生出的十个王，正如达尼尔所解释的那样。他说：\Quote{我注视那兽，就是第四只；看，它后面有十个角，其中又生出一个小角，像嫩芽一样；它将拔除它前面的三个角。}这小角所指的正是那要复兴犹太人王国的假基督。而被它拔除的三个角，则指三个王，即埃及、利比亚和埃塞俄比亚的王，它将在战争中消灭它们；\Emph{并且}，当它击败所有人后，作为一个残暴的暴君，它将兴起对圣徒的磨难和迫害，自高自大攻击他们。\par

% structure: p0035 source=h2 target=section reason=relative-heading-rank; base=h2

\section{17}

你看，达尼尔如何向拿步高解释诸王国的统治；你看，他如何解释那像在各部分的形式；你已注意到，他如何预言性地指出四只兽从海中上来的含义。现在，我们有必要向你们详细揭示假基督所行的事，并尽我们所能，借助圣经和先知，向你们宣告他在全地的游荡和他无法无天的来临。\par

% structure: p0037 source=h2 target=section reason=relative-heading-rank; base=h2

\section{18}

正如主耶稣基督从圣洁无玷的童贞女取得肉身，居住在我们中间，并取了犹大支派，由此而出，圣经在雅各伯的祝福中宣告了他的王族谱系，当时雅各伯对他的儿子说：\Quote{犹大，你的兄弟要赞美你；你的手要压住你仇敌的颈项；你父亲的儿子要向你俯伏。犹大是幼狮；我儿，你从嫩芽上升；他蹲伏，他如雄狮、如母狮卧下，谁能唤醒他？权杖不离犹大，立法者不离他两腿之间，直到那为他所贮藏的来到，他是万民所期待的。}\BibleRef{创 49:8}请注意雅各伯对犹大所说的这些话，都在主身上应验了。同样，族长关于假基督也这样表达。因此，正如他对犹大预言，同样也对他的儿子丹预言。因为犹大是他的第四子，而丹是他的第七子。那么，他对丹说了什么呢？\Quote{让丹成为路旁的蛇，咬伤马蹄？}\BibleRef{创 49:17}那蛇不就是起初的欺骗者吗？就是在创世纪中得名的，他欺骗了厄娃，咬伤了亚当的脚跟？\par

% structure: p0039 source=h2 target=section reason=relative-heading-rank; base=h2

\section{19}

既然现在我们必须更详细地证明所主张的，我们就不会畏缩。因为可以肯定，他注定要从丹支派而生，并像一个王子般的暴君、可怕的审判者和控告者那样敌对，正如先知作证说：\Quote{丹要审判他的百姓，如同以色列中的一个支派。}\BibleRef{创 49:16}但有人可能会说，这话是指三松，他出自丹支派，审判了百姓二十年。然而，那在三松身上只是部分应验，而在假基督身上将完全应验。因为耶肋米亚也这样说：\Quote{从丹我们听到他战马的疾驰声；在他战马的嘶鸣声中，整个大地都震动。}\BibleRef{耶 8:16}此外，梅瑟说：\Quote{丹是幼狮，他要从巴商跃出。}\BibleRef{申 33:22}为使没有人误以为这话是指救主，请注意这一点。他说：\Quote{丹是幼狮}；通过这样称呼丹支派为控告者注定要出的支派，他使所谈的事十分清楚。因为正如基督出自犹大支派，同样假基督将出自丹支派。正如我们的主和救主耶稣基督、天主子，在预言中因他的王权和荣耀被称为\Emph{狮子}，同样，圣经也预言性地称控告者为狮子，因他的暴政和武力。\par

% structure: p0041 source=h2 target=section reason=relative-heading-rank; base=h2

\section{20}

因为那个欺骗者在各方面都试图使自己显得像天主子。基督是狮子，假基督也是狮子。基督是天上和地上万物的君王，假基督也将是地上的君王。救主显现如羔羊，他也要显现如羔羊，但内心是豺狼。救主受了割损，他也要显现在割损中。救主派遣宗徒到万民中，他也要派遣假宗徒。基督聚集了四散的羊群，他也要聚集四散的希伯来人。基督将尊贵而赋予生命的十字架赐予相信他的人，他也要给出自己的记号。基督以人的形像显现，他也要以人的形像出现。基督从希伯来人中复活，他也要从犹太人中出现。基督显示祂的肉身如同圣殿，并在第三日重建；他也要在耶路撒冷重建石头的圣殿。那些由他编造的欺骗，对于留心听我们的人，将从接下来所陈述的变得十分明白。\par

% structure: p0043 source=h2 target=section reason=relative-heading-rank; base=h2

\section{21}

因为通过圣经，我们得知基督和救主有两次来临。第一次是在肉身中，处于卑微，因为祂显现于贫贱的境遇。因此，祂的第二次来临被宣告要在荣耀中；因为祂将从天上带着权能、天使和祂圣父的荣耀而来。祂的第一次来临有\Person{洗者若翰}作为前驱；而祂第二次来临（祂要在荣耀中降临）将显出\Person{厄诺客}、\Person{厄里亚}和\Person{神视者若望}。看，主对人的仁慈何等深厚；即使在最后的时期，祂仍显示对凡人的眷顾，怜悯他们。因为祂甚至在那时也不会没有先知就撇下我们，而要派遣他们到我们这里，为教导我们、保证我们，并使我们留心那对头的来临，正如祂从前在这位\Person{达尼尔}中所暗示的。因为他说：\BibleQuote{我要订立一七的盟约，在那一七中间，我的祭献和奠祭将被废除。}因为藉着一七，他指出在末世要显出的那七年。那一七的一半，两位先知连同若望将用于向全世界宣告假基督的来临，也就是说，穿着苦衣\BibleQuote{一千二百六十天}，见：\BibleRef{默 11:3}。他们将施行奇迹异事，旨在使人羞愧悔改，即使藉着这些方法，因为他们极其无法无天和不虔敬。\BibleQuote{若有人要伤害他们，火将从他们口中出来，烧灭他们的仇敌。他们有权柄在假基督来临的日子闭塞天空，使雨不降下，又使水变成血，并随意用各样灾祸击打大地。}当他们宣告了这一切之后，他们将倒在刀下，被控告者所剪除。他们将完成他们的见证，正如\Person{达尼尔}也说的；因为他预见那从深渊上来的兽要与他们交战，即与\Person{厄诺客}、\Person{厄里亚}和\Person{若望}，并要战胜他们，杀死他们，因为他们拒绝将光荣归于控告者。那就是那冒出来的小角。而他心高气傲，在结束时开始自高自大，自视为天主，迫害圣徒并亵渎基督。\par

% structure: p0045 source=h2 target=section reason=relative-heading-rank; base=h2

\section{22}

但既然我们按着讨论的线索，不得不涉及那对头掌权的日子，就必须首先说明关乎他出生和成长的事；然后，我们必须如先前所说，将我们的论述转向解释这件事，即控告者和无法无天者之子要在各方面使自己相似我们的救主。如此，论证也使我们清楚明白这事。既然世界的救主为拯救人类，由无玷童贞玛利亚诞生，并以肉身的形式践踏了仇敌，行使了祂自己天主性的权能；同样，控告者也将从一个不洁的女子来到世上，却要由一个童贞女虚假地诞生。因为我们的天主曾藉着肉身与我们同在，那肉身正是祂为亚当和亚当所有后裔所造的血肉之躯，但无罪恶。然而控告者虽取了肉身，却只是外表；因为他怎能穿戴那并非他自己所造、反而每日与之争战的血肉呢？我亲爱的，依我看，他将这表象的肉身作为工具来使用。为此，他也要像一位神灵那样由童贞女出生，然后向其余的人显现为肉身。至于童女生子，我们只在全然圣洁的童贞者身上知道这事，她生了确实穿上肉身的救主。因为\Person{梅瑟}说：\BibleQuote{凡开胎的男性，都当称为圣归于主。}这绝不适于他；因为正如对头不会开胎，他也不会取真实的肉身，并像基督受割损那样受割损。又如同基督拣选了祂的宗徒，他也要为自己召集一整个人群，成为与他一样邪恶的门徒。\par

% structure: p0047 source=h2 target=section reason=relative-heading-rank; base=h2

\section{23}

此外，他尤其喜爱犹太民族。并且对这一切，他将施行奇事和可畏的异兆，那是假奇事而非真事，为欺骗他那不虔敬的同类。因为如果可能，他甚至会迷惑选民，使他们远离对基督的爱\BibleRef{玛 24:24}。但在他的最初步骤中，他将温和、可爱、安静、虔诚、和平、憎恶不义、厌恶贿赂、不容忍偶像崇拜；他说，他喜爱圣经，尊敬司铎，敬重长者，拒绝淫乱，憎恶奸淫，不听信谗言，不轻易起誓，善待旅客，怜悯穷人，富有同情心。然后他将施行奇事：洁净癞病人，使瘫痪者起来，驱逐魔鬼，宣告遥远之事如同眼前之事，复活死人，帮助寡妇，保护孤儿，爱众人，以爱调和那些争斗之人，并对他们说：\BibleQuote{不要让太阳在你愤怒中落下}\BibleRef{弗 4:26}。他不积蓄金子，不喜爱银子，也不寻求财富。\par

% structure: p0049 source=h2 target=section reason=relative-heading-rank; base=h2

\section{24}

所有这些，他将以败坏和欺诈的方式去做，目的在迷惑众人，使他们立他为王。因为当民族和支派看见他如此伟大的德行和权能时，他们都会同心合意地聚集起来，立他为王。而尤以希伯来民族为那暴君自己所喜爱，他们彼此说：在我们的世代中果真有这样一位如此良善和公义的人吗？这事尤其将发生在犹太民族身上，正如我先前所说；他们因以为将要在这样的权能中看见那位君王本人，就上前对他说：我们都信任你，并承认你在全地上是公义的；我们都盼望由你获得救恩；我们从你口中领受了公义且不腐败的审判。\par

% structure: p0051 source=h2 target=section reason=relative-heading-rank; base=h2

\section{25}

起初，那欺骗者、无法无天者，确将以狡诈的欺骗拒绝这样的荣耀；但众人坚持拥护他，并立他为王。此后他的内心将高傲，那从前温柔的人将变得凶暴，那追随爱德的人将变得毫无怜悯，那心里谦卑的人将变得傲慢和不近人情，那憎恨不义的人将迫害义人。然后，当他被抬举至他的王国时，他将发动战争；并会在愤怒中击打三位强大的国王——即埃及、利比亚和埃塞俄比亚的国王。此后，他将在耶路撒冷建造圣殿，很快又将之修复，并交给犹太人。然后，他将在心中高傲对抗每一个人；是的，他甚至会口出亵渎对抗天主，在他的欺骗中以为他将永远在地上为王；这可怜虫不知道，他的王国很快就要化为乌有，他将很快面对为他预备的火，以及所有信任他、服侍他的人。因为当达尼尔说：\Quote{我要在一周内订立我的盟约}，他指的是七年；半周是先知宣讲的时期，另一半周——即三年半——假基督将在地上为王。此后，他的王国和荣耀将被夺去。看哪，你们这些爱慕天主的人，那些日子将兴起何等的灾难，自从世界奠基以来从未有过，将来也不会再有，唯独在那段日子。那时，那无法无天者心中高傲，将聚集他那些具有人形的恶魔，并憎恶那些召他进入王国的人，将玷污许多灵魂。\par

% structure: p0053 source=h2 target=section reason=relative-heading-rank; base=h2

\section{26}

因为他要从恶魔中指派首领治理他们。他将不再显得虔诚，而是全然并在一切事上变得严酷、严厉、暴躁、易怒、可怖、反复无常、令人畏惧、阴郁、可憎、可厌、野蛮、报复心强、不法。他致力于将全人类抛入灭亡的深渊，将增多假征兆。因为当所有人在他的表演中向他欢呼时，他将大声呼喊，使众人站立之处震动：\Quote{你们众民、部落、民族，要认识我的大能、权柄和我王国的力量。有哪位君王像我这样伟大？有哪位大神能与我相比？谁能抵挡我的权柄？}在围观者眼前，他将移动山脉离开原地，他将干脚行于海上，他将从天上降火，他将白昼变黑夜、黑夜变白昼，他将随心所欲地移动太阳；总之，在那些看见他的人面前，他将借着虚幻显现的权能，使地上和海中的一切元素都服从他。因为当他尚未显明自己是灭亡之子时，就已经挑起并激起公开的战争，直至战斗和杀戮，那么当他自己亲自来到时，人们将看见他的真面目，他将施展何等诡计、欺骗和迷惑，意图引诱所有人，使他们离开真理之路和天国之门呢？\par

% structure: p0055 source=h2 target=section reason=relative-heading-rank; base=h2

\section{27}

此后，天不再降下露水，云不再降下雨水，地也不再出产果实，海充满腥臭，河川干涸，海中的鱼死亡，人因饥渴而死亡；父亲抱着儿子，母亲抱着女儿，一同死去，没有人埋葬他们。整个地面充满腐烂尸体的臭气。海不再接纳河水，变得像淤泥，充满无限的臭气。然后，全地将有极大的瘟疫，也有无法安慰的哀哭、无量的哭泣、不止的哀恸。那时，人们将认为在他们之前死去的人是有福的，并对他们说：\Quote{打开你们的坟墓，收容我们这些可怜的人吧；打开你们的巢穴，接纳你们可怜的亲属和相识吧。你们有福，因为你们没有看见我们的日子。你们有福，因为你们不必目睹我们痛苦的生活，不必经历这无法医治的瘟疫，不必忍受这困住我们灵魂的窘迫。}\par

% structure: p0057 source=h2 target=section reason=relative-heading-rank; base=h2

\section{28}

然后，那可憎者将派遣他的命令到各个政府，同时通过恶魔和可见的人的手，说：\Quote{一位大能君王在地上兴起了；你们都来敬拜他；都来观看他王国的力量：因为，看哪，他将赐予你们谷物；他将赐予你们美酒、大量财富和高贵的尊荣。因为全地和海洋都听从他的命令。你们都来归向他。}由于食物匮乏，所有人都将去敬拜他；他将在他们的右手和额上盖上他的标记，使人不能用右手将尊贵的十字架记号盖在额上；但他的右手被束缚了。从那时起，他不能再封印他的任何一个肢体，而是依附于那欺骗者并服侍他；在他身上没有悔改。这样的人立刻丧失了天主和人的眷顾，而那欺骗者因他可憎的封印，仅给他们少许食物。他的封印在额上和右手上的数目是\Quote{六百六十六}。\BibleRef{默 13:18}关于这数目我有一点意见，虽不能确定；因为在这个数目中，当用文字表达时，已经发现许多名字。然而，我们或许可以说，这同一个封印的书写将给我们\Emph{我否认}一词。因为在那些日子里，借着那些他的仆役——即偶像崇拜者——那苦毒的仇敌曾使用\Emph{我否认}一词，当时无法无天者强迫基督的证人起誓：\Quote{否认你的天主，那被钉十字架者。}\par

% structure: p0059 source=h2 target=section reason=relative-heading-rank; base=h2

\section{29}

在憎恨一切善者的仇敌之时，这印将如此，其内容如下：\Quote{我否认天地之创造者，我否认圣洗，我否认我（先前）的侍奉，并依附于你，我信你。}因为这就是先知\Person{哈诺客}和\Person{厄里亚}将要宣讲的：\Quote{不要信那将要来到并被看见的仇敌；因为他是仇敌、败坏者和丧亡之子，并欺骗你们；因此他要杀死你们，并用刀击打他们。}看那仇敌的诡诈，知道那欺骗者的阴谋，他如何企图完全迷惑人灵。因为他要显示他的魔鬼光明如天使，并带来无数无形的军队。并且他在众人面前显扬自己被提到天上，伴有号角和响声，并那些用难以言表的赞歌迎接他的人的强大呼喊；黑暗的继承人自己如同光一般照耀，时而升到天上，时而带着极大光荣降到地上，又命令魔鬼如同天使，极其畏惧战栗地执行他的命令。那时他要派遣魔鬼的部队到山间、洞穴和地下的隐密处，追踪那些躲藏起来不被看见的人，并领他们来朝拜他。顺从者，他以他的印盖印；拒绝顺从他的人，他以无比的痛苦和最残酷的折磨与诡计消灭他们，这样的事从未有过，人耳未曾听闻，人眼未曾看见。\par

% structure: p0061 source=h2 target=section reason=relative-heading-rank; base=h2

\section{30}

那时，凡战胜暴君者，有福了。因为他们要被显扬为比先前的见证人更光荣更高超；因为先前的见证人只战胜了他的爪牙，但这些却推翻并征服了控告者本人，即丧亡之子。因此，我们的君王\Person{耶稣基督}将以何等的赞词和冠冕装饰他们！\par

% structure: p0063 source=h2 target=section reason=relative-heading-rank; base=h2

\section{31}

但让我们回到正题。当人们接受了那印，并且找不到食物和水时，他们\Emph{必将}以痛苦的声音走近他，说：\Quote{给我们吃和喝，因为我们都在饥饿和种种困苦中昏晕；并命令天给我们雨水，并驱走那些吃人的野兽。}那时那狡猾者将以绝对的无人性嘲笑他们，回答说：\Quote{天拒绝降雨，地不再出产果实；那么我能从哪里给你们食物呢？}那时，听了这欺骗者的话，这些可怜的人将意识到这就是邪恶的控告者，并将在痛苦中哀伤，痛哭流涕，用手打脸，撕裂头发，用指甲抓破脸颊，同时彼此说：\Quote{灾祸啊！痛苦的契约啊！欺骗的盟约啊！巨大的不幸啊！我们怎么被欺骗者蛊惑了！我们怎么与他联合了！我们怎么陷入他的罗网了！我们怎么被他可憎的网捕获了！我们怎么听了圣经却不明白！}因为确实，那些沉溺于生活琐事和此世情欲的人，那时将容易归向控告者，并被他盖印。\par

% structure: p0065 source=h2 target=section reason=relative-heading-rank; base=h2

\section{32}

但许多聆听神圣圣经、手中有圣经、并以理解牢记在心的人，将逃脱他的欺骗。因为他们将清晰地看穿他阴险的外表和欺骗性的伪装，并逃离他的手，往山上去，躲藏在地下的洞穴中；他们将流泪并以痛悔的心寻求人类之友；祂必将他们从他的罗网中救出，并以祂的右手拯救那些以相称和正义的方式向祂恳求的人脱离他的圈套。\par

% structure: p0067 source=h2 target=section reason=relative-heading-rank; base=h2

\section{33}

你们看见了圣人们那时将以何种的斋戒和祈祷来修炼自己。也要注意，那将要临到城乡之人的季节和时期是多么艰难。那时他们要从东被带到西；并从西来到东，大大哭泣，大声哀号。当白昼开始破晓时，他们渴望黑夜，以便从劳苦中得到安息；而当黑夜笼罩他们时，由于持续的地震和空中的暴风雨，他们甚至渴望看见白昼的光，并寻求如何以后遇见惨死。那时整个大地将哀痛痛苦的生命，海洋和空气同样也要哀痛；太阳也要哀哭；野兽和飞鸟也要哀哭；山岳和丘陵，以及平原的树木，都要因人类而哀哭，因为所有人都离弃了圣洁的天主，听从了欺骗者，并接受了那个可憎者——天主的仇敌——的标记，而不是救主赋予生命的十字架。\par

% structure: p0069 source=h2 target=section reason=relative-heading-rank; base=h2

\section{34}

教会也要大大哀哭，因为无论\Quote{祭献和香}都不再被进行，也没有中悦天主的敬礼；但教会的圣所将成为像园中看守的茅舍，圣体圣血在那日子不再显现。天主的公共敬礼将被熄灭，圣咏歌唱止息，听不到圣经宣读；但人将有黑暗，哀号加哀号，灾祸加灾祸。那时金银被抛到街上，无人捡拾；但一切都被视为可憎。因为所有人都急于逃匿，却无处可逃避仇敌的灾祸；但由于他们身上带有他的标记，他们将轻易被认出并被宣告属于他。外面有恐惧，里面有战栗，不分昼夜。街道和房屋里有尸体；街道和房屋里有饥渴；街道上有骚乱，家中有哀哭。面容的美貌枯干，因为他们的形貌如同死人；妇女的美貌消退，所有男子的欲望消失。\par

% structure: p0071 source=h2 target=section reason=relative-heading-rank; base=h2

\section{35}

尽管如此，即使在那时，仁慈而良善的天主也不会使人性完全失去安慰；相反，祂将缩短那些日子和三年半的时期，并为了那些隐藏在深山和洞穴中的余民而减少那些时日，免得所有那些圣人的队伍完全灭绝。但这些日子将迅速过去；欺骗者和假基督的王国将被迅速除去。然后，最终，在一眨眼之间，这世界的形态将要消逝，人类的能力归于乌有，所有这些可见之物都将被毁灭。\par

% structure: p0073 source=h2 target=section reason=relative-heading-rank; base=h2

\section{36}

因此，亲爱的诸位，既然我们先前所说的事还在将来，当那一周被分割，招致荒凉的可憎之物已经兴起，主的先驱者已经完成了他们应尽的历程，最终整个世界达到终结时，除了我们从天上所盼望的主、救主耶稣基督、天主子——祂必带来火和一切公义的审判，针对那些拒绝相信祂的人——的显现之外，还剩下什么呢？因为主说：\BibleQuote{因为闪电从东方发出，直照到西方，人子的来临也要这样；无论在哪里有尸体，鹰也必聚集在那里。} \BibleRef{玛 24:27}因为十字架的标记将从东方直到西方升起，其光芒胜过太阳，宣布审判者的来临和显现，要按各人的行为报应各人。因为关于普遍的复活和圣人们的王国，达尼尔说：\BibleQuote{许多睡在尘土中的人要醒来，有些得永生，有些受羞辱和永久的轻蔑。}依撒意亚说：\BibleQuote{死人要复活，坟墓里的人要醒来，地上的人要欢欣。}我们的主说：\BibleQuote{在那一天，许多人要听见天主子声音，听见的人要活。}\par

% structure: p0075 source=h2 target=section reason=relative-heading-rank; base=h2

\section{37}

因为那时号角要吹响，唤醒那些从地下深处沉睡的人，义人和罪人一起。每个家族、语言、民族和部族都要在一眨眼之间复活；他们要站在地面上，以说不出的恐惧和战栗等候公义而可畏的审判者的来临。因为火河要像愤怒的海洋一样狂涌而来，烧毁山岳和丘陵，使海洋消失，像蜡一样用热力熔化大气。天上的星辰要坠落，太阳要变为黑暗，月亮要变为血。天要像书卷般卷起；整个大地要因其上所行的事而被烧尽，就是人们以败坏所行的，淫乱、奸淫、谎言、不洁、偶像崇拜、凶杀和战争。因为那时将会有新天新地。\par

% structure: p0077 source=h2 target=section reason=relative-heading-rank; base=h2

\section{38}

那时，圣洁的天使将奉命奔跑，聚集所有民族，就是那可怕的号角声从沉睡中唤醒的万民。在基督的审判台前，那些曾经是君王、统治者、大司祭和司铎的人要站立，并要为他们治理的账目和他们的羊群——无论谁因疏忽而丢失了羊群中的一只羊——而交账。那时，那些不满足于自己的粮饷反而欺压寡妇、孤儿和乞丐的士兵要被带上来。那时，收税的人要被控告，他们掠夺穷人超过应得之数，并制造假金以欺骗贫困者，在田地、房屋和教会中。那时，好色之徒要羞耻地起来，他们没有保持床榻的纯洁，反而被各种肉体的美色所诱惑，随从自己的私欲而行。那时，那些没有持守主的爱的人要沉默阴郁地起来，因为他们轻视了救主那轻省的诫命，即\BibleQuote{你应爱近人如你自己。}那时，持有不公的天平、不公的法码和量器以及干量器的人也要哭泣，他们正等候公义的审判者。\par

% structure: p0079 source=h2 target=section reason=relative-heading-rank; base=h2

\section{39}

我们何必详述那些在审判台前受审的人呢？那时，义人要如太阳般发光，而恶人则显得沉默阴郁。因为义人和恶人都要复活，成为不朽的：义人，为永远受尊荣，并品尝不朽的喜乐；恶人，为在审判中永远受罚。每个人都在思考，他将为自己的行为、善或恶，向公义的审判者作出怎样的回答。每个人的行为都要环绕他，无论他是善是恶。因为天上的权势要动摇，恐惧和战栗要吞噬万物，无论是天、地还是地底下的一切。万口都要公开承认祂，承认那来施行公义审判的祂，大能的天主、万物的创造者。那时，天使、宝座、权势、宰制、统治，以及革鲁宾和色辣芬——他们有着许多眼睛和六只翅膀——都要带着恐惧和惊讶前来，大声呼喊：\BibleQuote{圣、圣、圣，万军的上主，全能者；天地充满你的荣耀。}而万王之王、万主之主，不偏待人的审判者、向每人施行公义的法官，将出现在祂可畏而崇高的宝座上；所有的凡人都将带着极大的恐惧和战栗看见祂的面容，无论是义人还是罪人。\par

% structure: p0081 source=h2 target=section reason=relative-heading-rank; base=h2

\section{40}

那时，\OriginalTerm{丧亡之子}{son of perdition}，即控告者，与其魔鬼及其仆役，将由严厉而不容情的天使押解前来。他们将被投入永不熄灭的火、永不沉睡的虫和外面的黑暗中。因为希伯来人将看见他以人形显现，就如他昔日\Emph{降来}时，藉童贞圣母降生成人，并正如他们钉他在十字架上时所见。他将向他们展示手上的\Emph{钉痕}和脚上的钉痕，肋旁被长矛刺透的伤痕，荆棘编成的冠冕，以及他可敬的十字架。希伯来人将一次而永远地看见这一切，他们要哀号痛哭，正如先知所宣告：\BibleQuote{他们要瞻望他们所刺透的}；那时也没有人帮助他们或怜悯他们，因为他们不肯悔改，也没有离开邪恶的道路。这些人与魔鬼及控告者一同进入永罚。\par

% structure: p0083 source=h2 target=section reason=relative-heading-rank; base=h2

\section{41}

那时，他将聚集万民，正如神圣福音所惊人宣告的。因为福音作者玛窦，或者更确切地说，主自己，在福音中说什么？\BibleQuote{当人子在自己的光荣中，与众天使一同降来时，那时，他要坐在光荣的宝座上，万民都要聚集在他面前；他要把他们彼此分开，如同牧人分开绵羊和山羊一样：把绵羊放在自己的右边，山羊在左边。那时，君王要对那些在他右边的说：我父所祝福的，你们来吧！承受自创世以来给你们预备了的国度吧。} 你们来吧，先知们，你们曾为我的名被驱逐。你们来吧，圣祖们，你们在我来临之前就服从了我，并渴望我的国度。你们来吧，宗徒们，你们在我降生成人时，曾同我一起受苦，并在福音中\Emph{与我一同受苦}。你们来吧，殉道者，你们在暴君面前承认了我，忍受了许多痛苦和磨难。你们来吧，\OriginalTerm{圣统}{hierarchs}，你们日夜无瑕地为我举行神圣礼仪，每天奉献我可敬的身体和血的祭献。\par

% structure: p0085 source=h2 target=section reason=relative-heading-rank; base=h2

\section{42}

你们来吧，圣人们，你们曾在山间、洞穴和地窟中克苦己身，以节制、祈祷和贞洁尊崇我的名。你们来吧，贞女们，你们渴慕我的洞房，除我之外不爱任何新郎，你们以自己的见证和生活方式与我结为伉俪，我是不朽不坏的\OriginalTerm{新郎}{Bridegroom}。你们来吧，贫穷人和外乡人的朋友。你们来吧，持守我爱的，因为我是爱。你们来吧，拥有和平的，因为我就是和平。你们，我父所祝福的，来承受预备给你们的国度吧！你们不贪图财富，怜悯穷人，援助孤儿，帮助寡妇，给口渴者喝，喂养饥饿者，收留外乡人，给赤身者衣服穿，探望病人，安慰监禁者，帮助盲人，保持\OriginalTerm{信德的印记}{seal of the faith}完整无缺，在教会中聚会，聆听我的圣经，渴慕我的话，昼夜遵守我的法律，如同精兵与我一同忍受艰难，力求中悦我，你们天上的君王。你们来吧，承受自创世以来为你们预备的国度。看，我的国度已准备妥当；看，乐园已敞开；看，我的不朽已展现其美丽。你们都来承受自创世以来为你们预备的国度。\par

% structure: p0087 source=h2 target=section reason=relative-heading-rank; base=h2

\section{43}

那时，义人要回答，因这伟大而惊奇的事实而震惊：众天使不能面对面注视的那一位，竟称他们为朋友。他们要向他呼喊道：主，我们几时见了你饥饿而供养了你？夫子，我们几时见了你口渴而给你喝？\OriginalTerm{可畏者}{the Terrible One}，我们几时见了你赤身露体而给你衣服穿？不朽者，我们几时见了你作客而收留了你？\OriginalTerm{爱人者}{the Friend of man}，我们几时见你患病或在监里而来看望你？你是永生者。你像父一样无始，与圣神同永恒。你是从无中创造万有的那一位。你是天使的君王。你连深渊也战栗。你以光为衣。你造了我们，用泥土塑造了我们。你塑造了无形之物。整个大地在你面前逃避，我们怎能善待你的王权和统治？\par

% structure: p0089 source=h2 target=section reason=relative-heading-rank; base=h2

\section{44}

那时，\OriginalTerm{万王之王}{King of kings}要回答他们说：凡你们对我这些最小兄弟中的一个所做的，就是对我做的。凡你们收留了我对你们所说的那些人，给他们衣服穿，喂养他们，给他们喝，我是指那些属于我肢体的穷人——你们就是对我做的。但你们进入自创世以来为你们预备的国度吧；永远享受我天上的父和圣洁的\OriginalTerm{赋予生命的圣神}{quickening Spirit}所赐予你们的。那时，有哪张嘴能说出那些福乐呢？\BibleQuote{眼所未见，耳所未闻，人心所未想到的，就是天主为爱他的人所准备的。}\par

% structure: p0091 source=h2 target=section reason=relative-heading-rank; base=h2

\section{45}

你们已听过无尽的喜乐，听过不可动摇的国度，听过永无尽头的祝福盛宴。现在也来学习公义的审判者和慈悲的天主将以何等的愤怒和忿怒对左边的人所说的痛苦之言：\BibleQuote{你们可诅咒的，离开我，到那给魔鬼和他的使者预备的永火里去吧！你们为自己预备了这些；也来享受吧！你们可诅咒的，离开我，到外面的黑暗里，到那不灭的火里去吧，那是给魔鬼和他的使者预备的。我造了你们，你们却归附了别人。是我从母胎中引出你们，你们却弃绝了我。是我以命令的话用尘土塑造了你们，你们却归附了别人。是我养育了你们，你们却事奉了别人。我规定了地和海来维持你们的生命，划定你们生命的界限，你们却不听从我的诫命。我造了光为你们，使你们享受白昼，也造了黑夜使你们休息；你们却以邪言恶语惹怒我，藐视我，向情欲敞开了门。你们这些作恶的人，离开我去吧！我不认识你们，我不承认你们：你们使自己成了另一个主——即魔鬼——的工人。与他一同继承黑暗、不灭的火、不睡的虫和咬牙切齿吧。}\par

% structure: p0093 source=h2 target=section reason=relative-heading-rank; base=h2

\section{46}

\BibleQuote{因为我饿了，你们没有给我吃的；我渴了，你们没有给我喝的；我作客，你们没有收留我；我赤身露体，你们没有给我穿；我患病，你们没有来看顾我；我在监里，你们没有来探望我。我造了你们的耳朵，为叫你们听圣经；你们却把它们用于鬼魔的歌曲、竖琴和戏谑。我造了你们的眼睛，为叫你们看见我诫命的光并遵守它们；你们却召来了奸淫和放荡，让它们看向各样不洁。我预备了你们的嘴巴，为发出朝拜、赞美、圣咏和属灵诗歌，并恒常诵读；你们却使它适于辱骂、发誓和亵渎，你们坐着说邻居的坏话。我造了你们的手，为叫你们在祈祷和恳求中伸开；你们却伸手去抢劫、谋杀和彼此杀害。我规定了你们的脚，为在福音平安的预备中行走，无论在我圣人们的教会或家中；你们却教它们跑向奸淫、淫乱、戏院、舞蹈和骄傲之处。}\par

% structure: p0095 source=h2 target=section reason=relative-heading-rank; base=h2

\section{47}

\Quote{最后，集会解散，今生的景象止息：它的欺骗和表象都过去了。你要依附我，天上、地上和地下的万膝都向我跪拜。因为所有疏忽、没有以善行显示怜悯的人，除了不灭的火之外，一无所有。因为我是人类的友人，然而也是公义的审判者。我要按功过报偿；我要按各人的辛劳给予酬报；我要按各人的争战回报。我渴望怜惜，但看不到你的器皿里有油。我渴望怜悯，但你们一生全无怜悯。我渴望慈悲，但你们的灯因心硬而昏暗。离开我去吧！因为对没有显示怜悯的人，审判是没有怜悯的。}\par

% structure: p0097 source=h2 target=section reason=relative-heading-rank; base=h2

\section{48}

\Quote{那时，他们也要对那不讲情面的可畏审判者回答说：\InnerQuote{主啊，我们什么时候见你饿了，或渴了，或作客，或赤身露体，或患病，或在监里，没有服事你呢？主啊，你不认识我们吗？你造了我们，你塑造了我们，你用四元素造了我们，你赐给我们灵魂和心灵。我们信了你，我们领受了你的印记，我们获得了你的圣洗，我们承认你是天主，我们知道你是造物主，我们因你行了异能，我们藉你驱逐了魔鬼，我们为你克苦了肉体，我们为你保持了童贞，我们为你实践了贞洁，我们为你在地上成了客旅；你却说：我不认识你们，离开我去吧！}那时，他要回答他们说：\InnerQuote{你们承认我是主，却没有遵守我的话。你们有了我十字架的印记，却因心硬而涂抹了它。你们获得了我的圣洗，却没有遵守我的诫命。你们使身体服从于童贞，却没有持守怜悯，没有从心中除去对弟兄的仇恨。因为凡称呼我主啊，主啊的，不一定都得救，唯有遵行我旨意的人。这些人要进入永罚，而义人却要进入永生。}}\par

% structure: p0099 source=h2 target=section reason=relative-heading-rank; base=h2

\section{49}

\BibleQuote{你要忠信至死，我必赐给你生命的华冠。}\par

你们已听见了，亲爱的，主的回答；你们学会了审判者的判决；你们明白了何等可怕的审判等待着我们，以及何样的日子和时辰摆在我们面前。因此，让我们每天思量这些；让我们日夜思量，在家中、路上、和教堂里，好使我们在那可畏而不偏袒的审判中不至被定为有罪、卑贱、愁苦，而是以纯洁的行为、生命、交谈和宣认站立；这样，慈悲而善良的天主也会对我们说：\BibleQuote{你的信德救了你，平安去吧；} 并且又说：\BibleQuote{好！善良忠信的仆人，你在小事上忠信，我必委派你管理许多大事：进入你主人的喜乐吧。} 愿我们也能获得那喜乐，因我们的主耶稣基督的恩宠和慈爱，愿光荣、尊威和朝拜归于他，连同他那无始的父，以及他的圣、慈善和赋予生命的圣神，从现在直到永远，世世代代。阿们。\par

\SourceRecord{Translated by J.H. MacMahon.；Ante-Nicene Fathers；Vol. 5.；Edited by Alexander Roberts, James Donaldson, and A. Cleveland Coxe.；Buffalo, NY: Christian Literature Publishing Co.,；1886.；<http://www.newadvent.org/fathers/0504.htm>.}
